% --- Archon physics formalization source begin ---
% archon:chemistry
% archon:covers IChO2026Problems/problem_icho_2026_t7_a3.lean
% archon:problem-id icho_2026_t7
% archon:part-id T7-A3
% archon:previous-part icho_2026_t7_a1
% archon:previous-part-policy derive_in_answer_blind_run_or_use_problem_stated_fallback

\chapter{Icho Chemistry problem icho\_2026\_t7\_a3}
\label{ch:IChO2026Problems_problem_icho_2026_t7_a3}

\paragraph{Problem source.}
\#\# Shared official problem context
T7. Nitrogen Fixation 7\% The Haber–Bosch process is the principal industrial method for ammonia production. A simplified diagram of an ammonia synthesis plant is shown below. Fig. 1. M1 - Mixture 1, M2 - Mixture 2, Z −CO2 scrubber, CLR - cooler In the following calculations, assume that all reactions in Fig. 1, except NH3 formation, are quantitative. 7.1 Tick the gases contained in M1 and M2; and tick the correct relationship between 𝑥and 𝑦. An ammonia production complex with a capacity of 660, 000 tons per year was commissioned at the Navoiazot plant in Uzbekistan, which operates with a 97.0\% of overall yield. 7.2 Calculate the mass of CH4 (tons) required annually to produce 660, 000 tons of ammonia, if the Navoiazot operates according to the Fig. 1. Consider a model of a real reaction system, which includes recirculation of reagents. An open system operates in a cyclic mode. Each cycle consists of four steps: 1. A stoichiometric N2 : H2 = 1 : 3 mixture with a total amount of n0 mol is fed into the reactor after each cycle. 2. The reaction proceeds with a constant yield of 𝜂= 0.150. 3. Ammonia is separated by liquification. 4. Unreacted gases are returned into the reactor. The cycle is repeated until the required overall yield is reached. Fig. 2. M2. – Mixture 2, CLR - cooler.

\#\# Current subquestion T7-A3
a) Calculate the amount of nitrogen, present in the system after the 58th cycle (before the addition of 59th portion of mixture), if 𝑛0 = 4 mol. Provide 4 decimal places. Hint: a formula for the sum of the terms of a geometric progression is needed. b) Calculate the number of cycles needed to increase the overall yield from 15.0\% to 97.0\%.

\paragraph{Shared problem context.}
T7. Nitrogen Fixation 7\% The Haber–Bosch process is the principal industrial method for ammonia production. A simplified diagram of an ammonia synthesis plant is shown below. Fig. 1. M1 - Mixture 1, M2 - Mixture 2, Z −CO2 scrubber, CLR - cooler In the following calculations, assume that all reactions in Fig. 1, except NH3 formation, are quantitative. 7.1 Tick the gases contained in M1 and M2; and tick the correct relationship between 𝑥and 𝑦. An ammonia production complex with a capacity of 660, 000 tons per year was commissioned at the Navoiazot plant in Uzbekistan, which operates with a 97.0\% of overall yield. 7.2 Calculate the mass of CH4 (tons) required annually to produce 660, 000 tons of ammonia, if the Navoiazot operates according to the Fig. 1. Consider a model of a real reaction system, which includes recirculation of reagents. An open system operates in a cyclic mode. Each cycle consists of four steps: 1. A stoichiometric N2 : H2 = 1 : 3 mixture with a total amount of n0 mol is fed into the reactor after each cycle. 2. The reaction proceeds with a constant yield of 𝜂= 0.150. 3. Ammonia is separated by liquification. 4. Unreacted gases are returned into the reactor. The cycle is repeated until the required overall yield is reached. Fig. 2. M2. – Mixture 2, CLR - cooler.

\paragraph{Current subquestion.}
a) Calculate the amount of nitrogen, present in the system after the 58th cycle (before the addition of 59th portion of mixture), if 𝑛0 = 4 mol. Provide 4 decimal places. Hint: a formula for the sum of the terms of a geometric progression is needed. b) Calculate the number of cycles needed to increase the overall yield from 15.0\% to 97.0\%.

\paragraph{Figure/image paths.}
\begin{itemize}
\item icho\_2026\_source/image/T7\_page-2.png
\item icho\_2026\_source/image/T7\_page-1.png
\end{itemize}

\paragraph{Reusable previous-part conclusions.}
\begin{itemize}
\item Source T7-A1. Question: Tick the gases contained in M1 and M2; and tick the correct relationship between 𝑥and 𝑦. Policy: derive\_in\_answer\_blind\_run\_or\_use\_problem\_stated\_fallback
\end{itemize}

\paragraph{Formalization target.}
create a compiling Lean file with sorry bodies at `IChO2026Problems/problem\_icho\_2026\_t7\_a3.lean`.
The Lean declarations must preserve every mathematical object, quantifier, hypothesis, side condition, approximation/error guarantee, and requested conclusion from the source statement.
Use Mathlib/Physlib/CRNT names found through LeanExplore where available. Prefer the configured benchmark Base library; introduce a local definition only when the source genuinely requires an object that the environment does not expose.
This is an answer-blind solve-phase task. Derive a candidate result only from the problem-side evidence above; do not assume a target value, select a tolerance around a desired value, or encode a desired result into a candidate domain.
For numerical questions, define the raw end-to-end quantity and apply an explicit source-derived rounding rule. For identification questions, characterize or prove uniqueness before recording the derived candidate.

\begin{theorem}[Icho Chemistry formalization target]
\label{thm:physics:icho_2026_t7_a3:target}
The assigned autoformalize agent should translate this IChO chemistry problem into Lean declarations in the covered file, with theorem and lemma proof bodies written as `by sorry`.
\end{theorem}
\begin{proof}
This is an autoformalization task, not a proof task. Produce faithful statements that can later be proved without weakening the source contract.
\end{proof}
% --- Archon physics formalization source end ---
